% GENERATED FILE. Edit canonical lessons, then run npm run generate:books.
% canonical-source-hash: fnv1a64:929cb9e7fc132114
% canonical-lessons: TA-C247-kuruvi, TA-C247-kaluku, TA-C247-naykkutti, TA-C247-anil, TA-C247-pulu

\chapter{Sparrow, Eagle, Puppy, Squirrel, Worm}
\label{ch:ta-sparrow-eagle-puppy-squirrel-worm}
\hlchaptermodality{247}

\begin{chapteropening}
\textbf{By the end of this chapter:} \emph{I can say a sparrow, an eagle, a puppy, a squirrel and a worm.}

Complete the last lesson of chapter 247: I can say a sparrow, an eagle, a puppy, a squirrel and a worm.
\end{chapteropening}

\section[\texorpdfstring{\ta{குருவி} (kuruvi)}{குருவி (kuruvi)}]{\ta{குருவி} (kuruvi) — a sparrow}

\label{lesson:TA-C247-kuruvi}



\begin{quote}
Before the new one: say the Tamil for a habit, then the Tamil for tradition.
\end{quote}

\subsection*{You'll want to know: \ta{குருவி}}
\textbf{\ta{குருவி}} — \emph{kuruvi} — \textquotedblleft{}a sparrow\textquotedblright{}.

\begin{grammarlens}[title={What you've built}]
One more word for this chapter.
\end{grammarlens}

\subsection*{Your turn}
\begin{itemize}
\raggedright
  \item \emph{Say it:} \emph{kuruvi}
  \item \emph{Say it:} \emph{kuruvi}, once more
  \item \emph{Say it:} say \emph{kuruvi}
  \item \emph{Recall:} say the Tamil for a habit, then the Tamil for tradition, then say \emph{kuruvi} again
\end{itemize}

\begin{tcolorbox}[breakable,colback=teal!4,colframe=teal!35!black,title={Before you move on}]
What does \ta{குருவி} mean? (\textquotedblleft{}A sparrow\textquotedblright{}.) Say it once more.
\end{tcolorbox}
\section[\texorpdfstring{\ta{கழுகு} (kaḻuku)}{கழுகு (kaḻuku)}]{\ta{கழுகு} (kaḻuku) — an eagle}

\label{lesson:TA-C247-kaluku}



\begin{quote}
Before the new one: say the Tamil for tradition, then the Tamil for a sparrow.
\end{quote}

\subsection*{You'll want to know: \ta{கழுகு}}
\textbf{\ta{கழுகு}} — \emph{kaḻuku} — \textquotedblleft{}an eagle\textquotedblright{}.

\begin{grammarlens}[title={What you've built}]
One more word for this chapter.
\end{grammarlens}

\subsection*{Your turn}
\begin{itemize}
\raggedright
  \item \emph{Say it:} \emph{kaḻuku}
  \item \emph{Say it:} \emph{kaḻuku}, once more
  \item \emph{Say it:} say \emph{kaḻuku}
  \item \emph{Recall:} say the Tamil for tradition, then the Tamil for a sparrow, then say \emph{kaḻuku} again
\end{itemize}

\begin{tcolorbox}[breakable,colback=teal!4,colframe=teal!35!black,title={Before you move on}]
What does \ta{கழுகு} mean? (\textquotedblleft{}An eagle\textquotedblright{}.) Say it once more.
\end{tcolorbox}
\section[\texorpdfstring{\ta{நாய்க்குட்டி} (nāykkuṭṭi)}{நாய்க்குட்டி (nāykkuṭṭi)}]{\ta{நாய்க்குட்டி} (nāykkuṭṭi) — a puppy}

\label{lesson:TA-C247-naykkutti}



\begin{quote}
Before the new one: say the Tamil for a sparrow, then the Tamil for an eagle.
\end{quote}

\subsection*{You'll want to know: \ta{நாய்க்குட்டி}}
\textbf{\ta{நாய்க்குட்டி}} — \emph{nāykkuṭṭi} — \textquotedblleft{}a puppy\textquotedblright{}.

\begin{grammarlens}[title={What you've built}]
One more word for this chapter.
\end{grammarlens}

\subsection*{Your turn}
\begin{itemize}
\raggedright
  \item \emph{Say it:} \emph{nāykkuṭṭi}
  \item \emph{Say it:} \emph{nāykkuṭṭi}, once more
  \item \emph{Say it:} say \emph{nāykkuṭṭi}
  \item \emph{Recall:} say the Tamil for a sparrow, then the Tamil for an eagle, then say \emph{nāykkuṭṭi} again
\end{itemize}

\begin{tcolorbox}[breakable,colback=teal!4,colframe=teal!35!black,title={Before you move on}]
What does \ta{நாய்க்குட்டி} mean? (\textquotedblleft{}A puppy\textquotedblright{}.) Say it once more.
\end{tcolorbox}
\section[\texorpdfstring{\ta{அணில்} (aṇil)}{அணில் (aṇil)}]{\ta{அணில்} (aṇil) — a squirrel}

\label{lesson:TA-C247-anil}



\begin{quote}
Before the new one: say the Tamil for an eagle, then the Tamil for a puppy.
\end{quote}

\subsection*{You'll want to know: \ta{அணில்}}
\textbf{\ta{அணில்}} — \emph{aṇil} — \textquotedblleft{}a squirrel\textquotedblright{}.

\begin{grammarlens}[title={What you've built}]
One more word for this chapter.
\end{grammarlens}

\subsection*{Your turn}
\begin{itemize}
\raggedright
  \item \emph{Say it:} \emph{aṇil}
  \item \emph{Say it:} \emph{aṇil}, once more
  \item \emph{Say it:} say \emph{aṇil}
  \item \emph{Recall:} say the Tamil for an eagle, then the Tamil for a puppy, then say \emph{aṇil} again
\end{itemize}

\begin{tcolorbox}[breakable,colback=teal!4,colframe=teal!35!black,title={Before you move on}]
What does \ta{அணில்} mean? (\textquotedblleft{}A squirrel\textquotedblright{}.) Say it once more.
\end{tcolorbox}
\section[\texorpdfstring{\ta{புழு} (puḻu)}{புழு (puḻu)}]{\ta{புழு} (puḻu) — a worm}

\label{lesson:TA-C247-pulu}



\begin{quote}
Before the new one: say the Tamil for a puppy, then the Tamil for a squirrel.
\end{quote}

\subsection*{You'll want to know: \ta{புழு}}
\textbf{\ta{புழு}} — \emph{puḻu} — \textquotedblleft{}a worm\textquotedblright{}.

\begin{grammarlens}[title={What you've built}]
One more word for this chapter.
\end{grammarlens}

\subsection*{Your turn}
\begin{itemize}
\raggedright
  \item \emph{Say it:} \emph{puḻu}
  \item \emph{Say it:} \emph{puḻu}, once more
  \item \emph{Say it:} say \emph{puḻu}
  \item \emph{Recall:} say the Tamil for a puppy, then the Tamil for a squirrel, then say \emph{puḻu} again
\end{itemize}

\begin{tcolorbox}[breakable,colback=teal!4,colframe=teal!35!black,title={Before you move on}]
What does \ta{புழு} mean? (\textquotedblleft{}A worm\textquotedblright{}.) Say it once more.
\end{tcolorbox}
